\ExamPart{Câu trắc nghiệm đúng sai}{2}{Trong mỗi ý a), b), c), d) ở mỗi câu, thí sinh chọn đúng hoặc sai.}

\NQuestion Cho phương trình $x^2-2(m+1)x+m^2+3=0$. Xét các mệnh đề sau:

\begin{SubQuestions}
\item Với mọi $m$, phương trình luôn là phương trình bậc hai.
\item Biệt thức của phương trình là $\Delta=8m-8$.
\item Khi $m=1$ thì phương trình có nghiệm kép.
\item Khi $m>1$ thì phương trình có hai nghiệm phân biệt.
\end{SubQuestions}

\NQuestion Trong mặt phẳng tọa độ $Oxy$, cho điểm $A(1;2)$ và đường thẳng $d:x-y+1=0$. Xét các mệnh đề sau:

\begin{SubQuestions}
\item Đường thẳng qua $A$ và song song với $d$ có phương trình $x-y+1=0$.
\item Đường thẳng qua $A$ và vuông góc với $d$ có một vectơ chỉ phương là $(1;1)$.
\item Khoảng cách từ $A$ đến $d$ bằng $\sqrt{2}$.
\item Hình chiếu vuông góc của $A$ lên trục $Ox$ là điểm $(1;0)$.
\end{SubQuestions}
